\scanpage{212}\mspage{157}

% MS 1–2
flèche de chgt de base, relative au \deleted{\unsure{foncteur|le foncteur} $f_!$,}\inserted{diagramme} commutatif de $\mathrm{Cat}$

% MS 3–5 upper square
\diagram{\begin{tikzcd}[ampersand replacement=\&]
 S' \arrow[d,"f"'] \& {}_{s\backslash}S' \arrow[l,"q"'] \arrow[d,"f'"] \\
 S \& e\simeq\{s\} \arrow[l,"p"'] \& (s\in\mathrm{Ob}\,S)
\end{tikzcd}}

% MS 6
savoir

% MS 7
\[
 p^*(f_!(\unsure{\xi|\zeta}'))\longrightarrow f'_!(q^*(\xi')
\]

% MS 8
est un isomorphisme (pour $\xi'\in\mathbb D(S')$).

% MS 9
Soit donc

% MS 10
\[
 \xi'\in\mathrm{Ob}\,\mathrm{HOT}_{\unsure{W|w}}(S'),\qquad
 \xi'=\mathrm{hot}_{S',W}(X')
\]

% MS 11
avec $X'$ dans $\mathrm{Cat}_{/S'}$. \unsure{Comme} \deleted{\unsure{le}}

% MS 12
\[
 f_!(\xi')=\mathrm{hot}_{S,W}(X'),
\]

% MS 13–19
et que $X'$ n'a \deleted{\unsure{t}} aucune raison d'être lisse sur $S$ (si on ne suppose pas $S'$ \inserted{lui-même} lisse sur $S$, auquel cas il pourrait être utile de \deleted{se \unsure{restreindre}} prendre $X'$ \deleted{\unsure{lui-même}} lisse sur $S'$, donc aussi sur $S$), pour calculer $p^*(f_!(\unsure{y|\xi}'))$, il n'est pas question de prendre la fibre $X'_s$, mais il faut prendre $ {}_{s\backslash}X'=\unsure{\Phi|\varPhi}_S(X')_{\unsure{-s|s}}$ :

% MS 20
\[
 p^*(f_!(\xi'))\simeq\mathrm{hot}_W({}_{s\backslash}X').
\]

% MS 21–24 lower left square pair
\diagram{\begin{tikzcd}[ampersand replacement=\&]
 X' \arrow[d] \& {}_{s\backslash}X' \arrow[l] \arrow[d] \\
 S' \arrow[d] \& {}_{s\backslash}S' \arrow[l,"q"'] \arrow[d] \\
 S \& {}_{s\backslash}S \arrow[l]
\end{tikzcd}}

% MS 21–27 beside and below squares
Considérons le diagramme de chgt. de base par $ {}_{s\backslash}S\longrightarrow S$ ci-contre, à carrés cartésiens. On voit d'abord que $q^*(\xi')$ n'est autre que \deleted{$\mathrm{hot}_{{}_{s\backslash}S',W}({}_{s\backslash}\unsure{S'|X'})$} la classe dans $\mathrm{HOT}_W({}_{s\backslash}S')$ de l'image inverse de $X'$ par le chgt de base $ {}_{s\backslash}S'\longrightarrow S'$. Pour l'assurer, il suffirait de prendre $X'$ lisse sur $S'$. \unsure{(ce qui|(celui qui}

% archive sheet number: 211
